\documentclass[11pt,a4paper]{article}

% ========================================
% PACKAGES
% ========================================
\usepackage[utf8]{inputenc}
\usepackage[T1]{fontenc}
\usepackage{lmodern}  % Better font support for larger sizes
\usepackage[english]{babel}
\usepackage[margin=2.5cm]{geometry}

% Mathematics
\usepackage{amsmath, amssymb, amsthm}
\usepackage{mathtools}

% Graphics and colors
\usepackage{graphicx}
\usepackage{xcolor}
\usepackage{tikz}
\usetikzlibrary{positioning, arrows.meta, shapes.geometric, calc}
\usepackage{pgfplots}
\pgfplotsset{compat=1.18}
\usepackage[most]{tcolorbox}

% Lists and formatting
\usepackage{enumitem}
\usepackage{parskip}
\usepackage{fancyhdr}
\usepackage{array}

% Code listings
\usepackage{listings}
\usepackage{algorithm}
\usepackage{algorithmic}

% Hyperlinks
\usepackage{hyperref}
\hypersetup{
    colorlinks=true,
    linkcolor=blue,
    urlcolor=blue,
    citecolor=blue
}

% ========================================
% THEOREM ENVIRONMENTS
% ========================================
\theoremstyle{definition}
\newtheorem{definition}{Definition}[section]
\newtheorem{example}{Example}[section]
\newtheorem{exercise}{Exercise}[section]

\theoremstyle{plain}
\newtheorem{theorem}{Theorem}[section]
\newtheorem{lemma}[theorem]{Lemma}
\newtheorem{proposition}[theorem]{Proposition}
\newtheorem{corollary}[theorem]{Corollary}

\theoremstyle{remark}
\newtheorem*{remark}{Remark}
\newtheorem*{observation}{Observation}

% ========================================
% CUSTOM COMMANDS
% ========================================
\newcommand{\R}{\mathbb{R}}
\newcommand{\N}{\mathbb{N}}
\newcommand{\Z}{\mathbb{Z}}
\newcommand{\Q}{\mathbb{Q}}
\newcommand{\C}{\mathbb{C}}

% ========================================
% TABLE OF CONTENTS DEPTH
% ========================================
\setcounter{tocdepth}{2} % Show only parts, sections, and subsections

% ========================================
% HEADER AND FOOTER
% ========================================
\setlength{\headheight}{14pt}
\pagestyle{fancy}
\fancyhf{}
\lhead{\leftmark}
\rhead{HCI}
\cfoot{\thepage}
\renewcommand{\sectionmark}[1]{\markboth{#1}{}}

% ========================================
% DOCUMENT INFORMATION
% ========================================
\title{\textbf{HCI - MultiModal Systems - Theory 2}\\
\large Artificial Intelligence}
\author{Jacopo Parretti}
\date{I Semester 2025-2026}

% ========================================
% DOCUMENT
% ========================================
\begin{document}

% Custom title page
\begin{titlepage}
    \begin{center}
        \vspace*{2cm}
        
        % University/Course info at top
        {\large\textsc{Artificial Intelligence}}
        
        \vspace{0.4cm}
        \textcolor{blue!60!black}{\rule{0.5\textwidth}{1pt}}
        
        \vspace{2.5cm}
        
        % Main title with modern spacing
        {\fontsize{36}{44}\selectfont\bfseries Human-Computer\\[0.5cm] Interaction}
        
        \vspace{1.5cm}
        
        % Subtitle with clean design
        \begin{tcolorbox}[
            colback=blue!3!white,
            colframe=blue!40!black,
            width=0.6\textwidth,
            arc=1mm,
            boxrule=0.8pt,
            left=15pt,
            right=15pt,
            top=10pt,
            bottom=10pt
        ]
            \centering
            {\large\textit{Lecture Notes}}\\[0.3cm]
            \textcolor{gray}{MultiModal Systems --- Theory 2}
        \end{tcolorbox}
        
        \vfill
        
        % Clean course details
        \begin{tcolorbox}[
            enhanced,
            colback=white,
            colframe=gray!30,
            width=0.5\textwidth,
            arc=0mm,
            boxrule=0.5pt,
            left=20pt,
            right=20pt,
            top=12pt,
            bottom=12pt,
            shadow={2mm}{-2mm}{0mm}{black!10}
        ]
            \centering
            {\small\textcolor{gray!70}{\textbf{ACADEMIC YEAR}}}\\[0.4cm]
            {\Large 2025--2026}
        \end{tcolorbox}
        
        \vspace{2cm}
        
        % Author
        {\LARGE\textsc{Jacopo Parretti}}
        
        \vspace{1.5cm}
        
    \end{center}
\end{titlepage}

\newpage
\tableofcontents
\newpage

\part{Foundations of Multimodal Interaction}

\section{Introduction to Intelligent Multimodal Interfaces}

\subsection{Course Objectives}

This course explores the fundamental theories and concepts of \textbf{Human-Computer Interaction (HCI)}, an interdisciplinary field that synthesizes knowledge from cognitive psychology, computer science, and design. The primary objectives are:

\begin{itemize}
    \item Understanding the theoretical foundations of human-computer interaction
    \item Developing practical skills in designing and implementing multimodal interfaces
    \item Exploring the integration of artificial intelligence techniques in interactive systems
    \item Analyzing nonverbal communication and its role in human-computer interaction
\end{itemize}

\subsection{Focus Areas}

The course places special emphasis on three interconnected dimensions:

\subsubsection{Technological Solutions}

Students will develop computer interfaces with a focus on both methodological and implementation aspects. The course emphasizes hands-on experience in building functional interactive systems, bridging the gap between theory and practice.

\subsubsection{Multimodal Interaction}

\begin{tcolorbox}[colback=blue!5!white,colframe=blue!75!black,title=\textbf{Multimodal Input and Output Modalities}]
Special attention is devoted to \textbf{multimodal solutions} that integrate multiple input and output modalities:
\begin{itemize}
    \item \textbf{Touch}: Tactile and haptic interaction
    \item \textbf{Vision}: Camera-based interaction and visual recognition
    \item \textbf{Natural Language}: Speech and text-based communication
    \item \textbf{Audio}: Sound-based interaction and auditory feedback
\end{itemize}
\end{tcolorbox}

\subsubsection{Intelligent Systems}

The course explores how \textbf{artificial intelligence techniques} can enhance interaction by:
\begin{itemize}
    \item Inferring user intentions from multimodal input
    \item Predicting expected interactions based on context and user behavior
    \item Adapting interface behavior to individual users
    \item Recognizing and responding to affective and social cues
\end{itemize}

\subsection{Course Program}

\subsubsection{Theoretical Component}

The theoretical component covers the following topics:

\begin{enumerate}
    \item \textbf{Introduction}: Course motivation, professional perspectives, open research issues, program overview, and examination methodology
    
    \item \textbf{Foundations of HCI}: Human factors in interface design, interaction design principles, usability evaluation, gaming, and gamification
    
    \item \textbf{Visual Interaction}: Camera calibration techniques, structure from motion, 3D reconstruction
    
    \item \textbf{Nonverbal Behavior in Communication}: 
    \begin{itemize}
        \item Types of nonverbal behavior: facial expressions, gestures, posture, eye gaze
        \item Data collection methods and protocols
        \item Tools and software for nonverbal behavior analysis
        \item Annotation tools (e.g., ELAN)
    \end{itemize}
    
    \item \textbf{Automated Analysis of Body Language}: Movement tracking, gesture recognition, facial expression analysis, and speech processing. Techniques for data capture, feature extraction, and automatic analysis
    
    \item \textbf{Social Artificial Intelligence}: Applications in social psychology, organizational psychology, and social robotics
    
    \item \textbf{Affective Computing}: Theories of emotion, emotion recognition systems, and applications in HCI
    
    \item \textbf{Multimodal Fusion}: Integration of multimodal nonverbal cues using fusion techniques (late fusion, early fusion)
\end{enumerate}

\subsubsection{Laboratory Component}

The laboratory sessions provide hands-on experience with state-of-the-art tools and techniques:

\begin{enumerate}
    \item \textbf{Deep Image Matching}: Python implementation of feature detection and matching algorithms
    
    \item \textbf{3D Model Reconstruction}: Structure from motion using Zephyr software
    
    \item \textbf{Camera Pose Estimation}: C\# implementation of Fiore's method for camera localization
    
    \item \textbf{3D Graphics}: Modeling and rendering in Unity game engine
    
    \item \textbf{Model-Based Augmented Reality}: Implementation of the complete AR pipeline integrating Python code and Unity
    
    \item \textbf{Advanced Topics}: Deep learning approaches to camera pose estimation and model recognition
\end{enumerate}
 
\newpage

\section{The Evolution of Human-Computer Interaction}

\subsection{Why Study HCI?}

The fundamental motivation for studying Human-Computer Interaction lies in understanding the \textbf{interaction space} between humans and machines. Effective interaction design must strike a balance between these two poles—neither overly machine-centric nor purely human-centric, but rather positioned at an optimal middle ground that accommodates both human capabilities and computational constraints.

The central challenge of HCI is to design interfaces that allow users to accomplish their goals efficiently and naturally, while leveraging the computational power and capabilities of modern computing systems.

\subsection{Historical Evolution of Interaction Paradigms}

The field of HCI has undergone several paradigm shifts, each bringing the interaction closer to natural human behavior.

\subsubsection{Command-Line Interfaces (Pre-1980s)}

\begin{tcolorbox}[colback=green!5!white,colframe=green!75!black,title=\textbf{Command-Line Interfaces (CLI)}]
Early human-computer interaction relied on \textbf{command-line interfaces (CLI)}, which required users to communicate with machines through text-based commands. This paradigm represented an interaction model heavily biased toward the machine:
\begin{itemize}
    \item \textbf{Expert-oriented}: CLIs assume users possess detailed knowledge of the system's internal organization and command syntax
    \item \textbf{Machine-centric}: The interface reflects the machine's architecture rather than human cognitive models
    \item \textbf{High learning curve}: Users must memorize commands, parameters, and system-specific conventions
    \item \textbf{Limited accessibility}: Generic users cannot interact effectively without substantial training
\end{itemize}
\end{tcolorbox}

While command-line interfaces persist today (particularly among expert users who value direct system control and automation capabilities), they exemplify an interaction paradigm positioned closer to the machine than to the human.

\subsubsection{Graphical User Interfaces (1980s–2000s)}

\begin{tcolorbox}[colback=orange!5!white,colframe=orange!75!black,title=\textbf{Graphical User Interfaces (GUIs)}]
The introduction of \textbf{Graphical User Interfaces (GUIs)} in the early 1980s marked a fundamental shift in interaction design:
\begin{itemize}
    \item \textbf{Direct manipulation}: Users interact with visual representations of objects through pointing devices (mouse) and keyboards
    \item \textbf{Visual metaphors}: Desktop metaphors, windows, icons, and menus provide intuitive representations of system functionality
    \item \textbf{Reduced cognitive load}: Users no longer need to memorize complex command syntax
    \item \textbf{Increased accessibility}: Non-expert users can accomplish tasks through exploration and recognition rather than recall
\end{itemize}
\end{tcolorbox}

This paradigm shift moved interaction significantly closer to human cognitive models, making computing accessible to a broader population.

\subsubsection{Pervasive and Ubiquitous Interaction (2000s–Present)}

\begin{tcolorbox}[colback=purple!5!white,colframe=purple!75!black,title=\textbf{Pervasive and Ubiquitous Computing}]
Over the past 15–20 years, a new paradigm has emerged: \textbf{pervasive interaction} (also known as ubiquitous computing). This paradigm is characterized by:
\begin{itemize}
    \item \textbf{Distributed computing}: Intelligence is distributed across multiple interconnected devices (smartphones, wearables, IoT sensors, smart home devices)
    \item \textbf{Invisible interfaces}: The traditional notion of "the machine" becomes diffuse—computation is embedded in the environment
    \item \textbf{Context awareness}: Systems adapt to user context, location, and activity
    \item \textbf{Natural interaction}: Voice, gesture, and multimodal input replace traditional keyboard and mouse
    \item \textbf{Smart environments}: Individual devices collaborate to create intelligent, responsive spaces
\end{itemize}
\end{tcolorbox}

In contrast to earlier paradigms where "the machine" was a discrete entity contained in a single case, pervasive computing dissolves the boundaries between user, environment, and computational infrastructure. This represents the closest approximation yet to natural human interaction patterns.

\subsection{Technological Foundations}

This course emphasizes the \textbf{technological aspects} that enable modern HCI systems. Implementing sophisticated, intelligent, and multimodal interactions requires substantial infrastructure:

\begin{itemize}
    \item \textbf{Network infrastructure}: High-bandwidth internet connectivity, satellite networks, and edge computing for distributed systems
    \item \textbf{Interaction devices}: Advanced input/output hardware including VR/AR headsets, depth cameras, haptic devices, and multimodal sensors
    \item \textbf{Computational power}: GPU acceleration, cloud computing resources, and specialized AI hardware for real-time processing
    \item \textbf{Software frameworks}: Machine learning libraries, computer vision toolkits, and multimodal fusion platforms
\end{itemize}

Understanding these technological foundations is essential for designing and implementing the next generation of human-computer interfaces. 



\section{Human Factors}

\textbf{Human factors} are the psychological and physiological characteristics of the human body that are relevant to the design of human-computer interfaces.
Perceptual aspects such as vision, hearing, touch, taste, and smell are relevant to the design of human-computer interfaces.
We will mainly focus on visual aspects.

\subsection{Brightness perception}
\textit{Brightness perception} is the ability to perceive the brightness of a light.

\begin{itemize}
    \item \textbf{Brightness} perception is subjective and depends on the individual.
    \item \textbf{Brightness} is conditioned by the quantity of light.
    \item \textbf{Brightness} is a function of the object brightness and the background brightness.
    \item \textbf{Visual acuity} is the ability to see details and increases with lighter scenes.
\end{itemize}

The \textit{flicker effect} perceived in large monitors is due to the refresh rate of the monitor.
A monitor with a refresh rate of 60 Hz refreshes the image 60 times per second.
If the image is not refreshed at this rate, the eye will perceive a \textit{flicker effect}.

\subsection{Abilities and limitations of vision processing}
The visual system compensates for motion (i.e., \textit{image stabilization}) and light (or color) variations.
Examples of overcompensations are demonstrated in the following optical illusions:

\subsection{Context awareness}
Context is important for the perception of the environment.
For instance, the perception of a face is different in a dark room than in a bright room. 
Limits of vision processing are also affected by the context.

While reading we follow a few steps:
\begin{itemize}
    \item perception of visual patterns,
    \item patterns are decoded using the internal representation of the 
    language,
    \item they are explained using syntactic and semantic 
    knowledge (and pragmatics)
\end{itemize}
Reading is performed by \textit{saccadic eye movement} and \textit{visual fixation}.
The shape of a single word is very important for the perception of the word.

\subsection{Cognitive aspects of vision processing}
Some of the cognitive aspects of vision processing are the following:
 \begin{itemize}
    \item Long-term memory and learning
    \item Problem solving
    \item Decision making
    \item Attention and problem dimensionality
\end{itemize}

\subsection{Human Memory}
We can differentiate between different \textit{sensorial buffers}:
\begin{itemize}
    \item \textit{iconic memory}
    \item \textit{echoic memory}
    \item \textit{haptic memory}
\end{itemize}
%sensorial buffer->attention -> short term memory -> long term memory
These sensorial buffers after being processed by the \textit{attention mechanism} are stored in the \textbf{short-term memory}.
The \textbf{short-term memory} is then processed by the \textbf{working memory} and stored in the \textbf{long-term memory} after 
enough repetitions.

\subsection{Short-term memory (STM)}

\begin{tcolorbox}[colback=yellow!5!white,colframe=yellow!75!black,title=\textbf{Miller's Law and Chunking}]
Information is remembered better when organized in chunks.
The chunk organization is called \textit{closure} and is classically summarized by \textit{Miller's Law}.
In general a user can remember 
%formula 7 plus or minus 2
\begin{math}
    7 \pm 2 \text{ chunks}
\end{math}
The user always expects a closure at the end of a task. For example, in early cash retrieval
machines users often forgot to take the card after getting the cash.
\end{tcolorbox}


\subsection{Inductive reasoning}
\textit{Inductive reasoning} describes the generalization of unseen examples from seen examples:
For example, if all cars I have seen have four wheels, then all cars have four wheels.
This is not fully reliable, as we can prove falsity but not absolute truth.
Humans are generally not able to use negative evidence effectively.

\section{Human and Virtual Reality}
\textbf{Virtual reality} is a technology that allows the user to interact with a virtual environment.
There are many issues about human perception that must be considered in depth when analyzing interaction in VR.
These issues are often neglected for classic \textbf{HCI} based on intentional actions and simple stimuli.
Some of the issues are the following:
\begin{itemize}
    \item Awareness, attention
    \item Information fusion
    \item Hierarchical processing
    \item Adaptation
    \item Psychophysics
\end{itemize}

\subsection{Human senses}

Humans have more than five senses.
Some of these senses include:
\begin{itemize}
    \item \textbf{Proprioception}: sense of body position—what is your
    body doing right now
    \item \textbf{Equilibrium}: balance
    \item \textbf{Acceleration}: perception of linear and angular changes
    \item \textbf{Nociception}: sense of pain
    \item \textbf{Thermoreception}: temperature
    \item \textbf{Satiety}
    \item \textbf{Thirst}
    \item \textbf{Micturition}
    \item \textbf{Chemoreception}: amount of CO$_2$ and Na$^+$ in blood
\end{itemize}

In virtual reality \textbf{vision} is the main channel, but other senses are also 
important and therefore must be coherent with the virtual environment.
\textbf{Sensing} is an active process that depends on previous inputs \textbf{not in a conscious way}.



\subsection{Adaptation}
\textbf{Adaptation} is the process of changing the perception of a stimulus over time.
For instance, the perceived loudness of motor noise in an aircraft or car decreases
within minutes.
The optical system of our eyes and the photoreceptor sensitivities
adapt to change perceived brightness.
Over long periods, perceptual training can lead to adaptation;
In military training simulations, sickness experienced by soldiers
appears to be less than expected, perhaps due to regular exposure.
The same appears to be true of experienced video game players. Those who have spent many hours and days in front of large screens
playing first-person shooter games apparently experience less
\textit{vection} when locomoting themselves in VR.
\textbf{Adaptation} therefore becomes a crucial factor for VR.


\subsection{Constraints for VR systems}
In VR systems we have to consider the following constraints due to human vision:
\begin{itemize}
    \item \textbf{Human frontal FOV}: 100$^\circ$--110$^\circ$ horizontal per eye and up to 200$^\circ$--220$^\circ$ with both eyes
    \item \textbf{Stereo overlap}: roughly 110$^\circ$--120$^\circ$
    \item \textbf{Foveal field}: approximately 60$^\circ$ horizontally and vertically where both eyes can see in focus
\end{itemize}

\subsection{Eye movements}
When designing AR/VR systems, we must take into account how human eyes move.
\begin{itemize}
    \item \textbf{Saccades}: brisk motions lasting less than 45 ms with rotations of about 90$^\circ$ per second that quickly relocate the \textit{fovea} so important features in a scene are sensed with the highest visual acuity, giving the illusion of high acuity over a large angular range. We typically have little awareness of these eye movements even though we can consciously control \textit{fixation targets}.
    \item \textbf{Vestibulo-ocular reflex (VOR)}: an involuntary movement of the eyes that counteracts head motion and is one of the most important motions to understand for VR.
    \item \textbf{Smooth pursuit}: slow eye movement (usually $<30^\circ$/s) used to track a moving target feature such as a car, tennis ball, or person walking by, reducing \textit{motion blur} on the retina.
    \item \textbf{Optokinetic reflex}: rapid movement of the eyes to track a fast object; the eyes rapidly and involuntarily choose features for tracking on the object.
    \item \textbf{Vergence stereopsis}: movement of the eyes to converge or diverge the pupils, providing important information about the distance of objects.
    \item \textbf{Microsaccades}: small, involuntary jerks (less than one degree) that trace out an erratic path and are believed to augment fixation control, reduce \textit{perceptual fading} due to adaptation, and improve visual acuity.
\end{itemize}


\subsection{Vergence and accommodation}
\textbf{Vergence} is the movement of the eyes to converge or diverge the pupils.
\textbf{Accommodation} is the adjustment of the eyes to focus on near objects.
\textbf{Binocular disparity} is the difference in the image seen by the two eyes.
\textbf{Retinal blur} is the blurring of the image seen by the eyes due to the finite size of the retina.

\begin{tcolorbox}[colback=red!5!white,colframe=red!75!black,title=\textbf{Vergence-Accommodation Mismatch in VR}]
In the real world, \textit{vergence} and \textit{accommodation} match, but
in VR, \textit{vergence} and \textit{accommodation} can mismatch.
This can cause discomfort:
\begin{itemize}
    \item \textbf{fatigue}
    \item \textbf{nausea}
    \item \textbf{eye strain}
    \item \textbf{motion sickness}
\end{itemize}
\textbf{Flicker fusion rate} is the frequency at which the eyes can perceive a flickering image.
It is important because it affects the perception of the image and can cause discomfort.
For instance, if the flicker fusion rate is too low, the eyes will perceive the image as flickering.

\textbf{Dynamic range} is the ratio of the brightest to the darkest part of the image.
Human vision has a dynamic range far higher than any available display technology.
\end{tcolorbox}


\subsection{Human Vision System}
The human vision system is the following:

\begin{table}[H]
    \centering
    \begin{tabular}{|l|>{\raggedright\arraybackslash}p{6.5cm}|}
    \hline
    \textbf{Visual acuity} & 20/20 is $\sim$1 arc min \\
    \hline
    \textbf{Field of view} & $\sim$200$^\circ$ monocular, $\sim$120$^\circ$ binocular, $\sim$135$^\circ$ vertical \\
    \hline
    \textbf{Resolution of eye} & $\sim$576 megapixels \\
    \hline
    \textbf{Temporal resolution} & $\sim$60 Hz (depends on contrast, luminance) \\
    \hline
    \textbf{Dynamic range} & instantaneous 6.5 f-stops, adapts to 46.5 f-stops \\
    \hline
    \textbf{Colour} & everything in CIE xy diagram \\
    \hline
    \textbf{Depth cues in 3D displays} & vergence, focus, (dis)comfort \\
    \hline
    \textbf{Accommodation range} & $\sim$8 cm to $\infty$, degrades with age \\
    \hline
    \end{tabular}
    \caption{Human Vision System}
\end{table}
Modern displays can't exactly match this level of performance. 
Let's see a comparison between the human eyes and the HTC Vive (2016) VR headset:

    \begin{table}[H]
    \centering
    \begin{tabular}{|l|c|c|}
    \hline
     & \textbf{Human Eyes} & \textbf{HTC Vive} \\
    \hline
    FOV & 200$^\circ$ $\times$ 135$^\circ$ & 110$^\circ$ $\times$ 110$^\circ$ \\
    \hline
    Stereo Overlap & 120$^\circ$ & 110$^\circ$ \\
    \hline
    Resolution & 30,000 $\times$ 20,000 & 2,160 $\times$ 1,200 \\
    \hline
    Pixels/inch & $>$2190 (100mm to screen) & 456 \\
    \hline
    Update & 60 Hz & 90 Hz \\
    \hline
    \end{tabular}
    \caption{Comparison between the human eyes and the HTC Vive (2016) VR headset}
\end{table}
    

The following are some of the 3D cues that are used to create the illusion of depth:
\begin{itemize}
    \item \textbf{Atmospheric artifacts}: fog, smoke, dust, etc.
    \item \textbf{Relative sizes}: objects that are closer are larger than objects that are further away.
    \item \textbf{Image blur}: objects that are further away are blurred.
    \item \textbf{Occlusion}: objects that are further away are occluded by objects that are closer.
    \item \textbf{Shadows}: objects that are further away have smaller shadows.
    \item \textbf{Texture}: objects that are further away have less texture.
    \item \textbf{Colour}: bluish objects are more distant.
    \item \textbf{Motion parallax}: Objects in relative motion provide depth information; closer objects appear to move faster than distant objects when the observer moves.
\end{itemize}


\subsection{Issues for VR}
The following are some of the issues for VR:
\begin{itemize}
    \item \textbf{Coherence of depth cues}: Achieving coherent depth cues across all visual elements is difficult.
    \item \textbf{Vergence-accommodation mismatch in VR headsets}: In conventional VR displays, vergence and accommodation are decoupled because the display is at a fixed focal distance while the eyes converge at different distances. This mismatch can cause discomfort and is a fundamental limitation of current VR technology.
    \item \textbf{Imperfect head tracking}: If there is significant latency, visual stimuli will not appear in the correct place at the expected time, causing spatial misalignment and potential motion sickness.
    \item \textbf{Limited tracking systems}: Many tracking systems monitor only head orientation, which has obvious drawbacks for accurate spatial interaction.

\end{itemize}


\subsection{Motion perception}
Motion detection in the human visual system depends on local processing at the \textit{retinal level}, which is then integrated at higher \textit{cortical levels} to create a coherent perception of motion.

\begin{itemize}
    \item \textbf{Optical flow}: The pattern of apparent motion of objects, surfaces, and edges in a visual scene caused by the relative motion between an observer and the scene. It represents the motion field projected onto the image plane and is fundamental to understanding how we perceive movement in both real and virtual environments.
    
    \item \textbf{Apparent motion in Virtual Reality}: VR systems exploit \textit{apparent motion} by rapidly displaying sequences of static images. The brain perceives continuous motion from discrete frames, similar to how \textit{motion pictures} and animation work. This illusion allows VR to create convincing motion experiences despite rendering discrete frames.
    
    \item \textbf{Stroboscopic apparent motion}: A visual phenomenon where discrete images presented in rapid succession are perceived as continuous motion. This is the fundamental principle behind animation, film, and VR displays, where the brain fills in the gaps between frames to create the perception of \textit{smooth motion}.
    
    \item \textbf{Beta and Phi effects}: Two related phenomena that explain apparent motion:
    \begin{itemize}
        \item \textbf{Beta effect}: The perception of motion between two static images when they are presented in rapid succession with appropriate timing and spatial separation.
        \item \textbf{Phi effect}: The perception of motion in a single object that appears to move between different positions, creating an illusion of continuous movement.
    \end{itemize}
    
    \item \textbf{Persistence of Vision}: A traditional but somewhat misleading explanation suggesting that images persist on the retina for approximately 100 ms after the stimulus is removed. However, modern research indicates that motion perception in film and VR is better explained by the \textit{phi phenomenon} and \textit{beta effect} rather than true \textit{persistence of vision}.
    
    \item \textbf{Motion perception at low frame rates}: Interestingly, motion can be perceived even at very low frame rates (as low as 2 fps), though it appears jerky and discontinuous. This observation challenges the traditional \textit{persistence of vision} explanation and demonstrates that motion perception is more complex than simple \textit{retinal persistence}.
\end{itemize}



\subsection{Issues in HMD}
The following are some of the issues in HMD:
\begin{itemize}
    \item \textbf{Reducing image display time}: Reducing the time each image is displayed can help minimize flickering.
    \item \textbf{Frame rate requirements}: Flickering appears if the frame rate is not high enough; frame rates greater than 90 fps are typically acceptable.
    \item \textbf{Fast pixel switching}: Fast pixel switching is needed to reduce flickering. OLED displays achieve switching times of less than 0.1 ms.
\end{itemize}

\subsection{Other important senses in VR}
In immersive environments \textbf{somatosensory cues} complement vision and audition by providing information about contact, pressure, body pose, and motion. The somatosensory system is a distributed network of receptors in the skin, muscles, tendons, and joints that continuously reports the state of the body. Key aspects for VR design include:
\begin{itemize}
    \item \textbf{Haptic sensations}: The skin contains specialized cells that respond to pressure, vibration, and temperature. Haptic feedback adds realism by matching virtual contact events with tactile effects such as vibration or force.
    
    \item \textbf{Tactile channel}: Tactile receptors do not have uniform density across the body. Areas such as the tongue and fingertips contain many more mechanoreceptors, allowing finer spatial discrimination of textures and edges.
    
    \item \textbf{Mechanoreceptors}: Specialized receptors that encode different aspects of touch (high-frequency vibration, stretch, low-frequency vibration, and steady pressure, respectively). Haptic devices stimulate these receptors indirectly through actuators embedded in controllers, gloves, or suits.
    
    \item \textbf{Proprioception (joint position sense)}: Sensors embedded in muscles, tendons, and joints report limb positions even with eyes closed. This ability lets users reach or touch virtual objects with confidence. Joints closer to the torso are generally sensed more accurately, and healthy adults localize hand position within $\sim$8 cm without visual feedback.
    
    \item \textbf{Kinaesthesia (joint motion sense)}: Movement is detected through dynamic responses of muscle spindles and tendon organs. Accurate motion cues help the user perceive forces and resistances generated by the virtual world.
\end{itemize}

\subsubsection{Two-point discrimination}
The \textit{two-point test} measures tactile acuity by determining the minimum separation at which two simultaneous touches are perceived as distinct. Lower values indicate higher spatial sensitivity. Typical thresholds are summarized in Table~\ref{tab:two-point}.

\begin{table}[H]
    \centering
    \begin{tabular}{|l|c|}
        \hline
        \textbf{Body site} & \textbf{Threshold distance (mm)} \\
        \hline
        Finger & $2$--$3$ \\
        Cheek & $6$ \\
        Nose & $7$ \\
        Palm & $10$ \\
        Forehead & $15$ \\
        Foot & $20$ \\
        Belly & $30$ \\
        Forearm & $35$ \\
        Upper arm & $39$ \\
        Back & $39$ \\
        Shoulder & $41$ \\
        Thigh & $42$ \\
        Calf & $45$ \\
        \hline
    \end{tabular}
    \caption{Approximate two-point discrimination thresholds across the body.}
    \label{tab:two-point}

\end{table}

\subsection{Proprioception and Kinesthesia: Advanced Aspects}

\subsubsection{Efference Copies and Motor Control}

The \textit{motor cortex}, which controls body movement, sends signals called \textbf{efference copies} to other parts of the brain to communicate which movements have been executed. This mechanism is analogous to having encoders on robot joints or wheels to indicate how much they have moved.

\textbf{Implication:} You cannot tickle yourself because your brain knows where your fingers are moving through these efference copies, which cancel out the expected sensory feedback.

\subsection{The Vestibular System}

The \textbf{vestibular system}, located in the inner ear alongside the \textit{cochlea}, provides crucial information about balance, spatial orientation, and motion:

\begin{itemize}
    \item \textbf{Utricle and saccule}: Together form the \textbf{otolithic system}, which measures linear acceleration. \textit{Otoliths} (calcium carbonate crystals) move over \textit{hair cells} to detect gravity and linear acceleration, providing information about head position relative to gravity.
    
    \item \textbf{Semicircular canals}: Three fluid-filled canals oriented in orthogonal planes measure angular acceleration (rotational head movements). Each canal is sensitive to rotation in its specific plane.
\end{itemize}

\subsection{Auditory Perception}

Similar to vision, auditory experiences are often surprising and based on adaptation, missing data, and assumptions filled by \textit{neural structures}. Many interesting aspects derive from \textit{psychophysical measurements}.

\subsubsection{Auditory Illusions}

\begin{itemize}
    \item \textbf{Precedence effect}: When two nearly identical sounds arrive at slightly different times, we perceive a single sound. This is fundamental for handling \textbf{reverberation} (delayed copies of sound due to reflections), allowing us to perceive a single sound source based on the first arrival (usually the strongest). An \textit{echo} is perceived if the temporal difference exceeds the "\textit{echo threshold}."
    
    \item \textbf{McGurk effect}: Demonstrates the power of multisensory integration by mixing visual and auditory cues. In experiments where a video of a person speaking is dubbed with mismatched audio (e.g., hearing "ba" while seeing "ga"), most subjects perceive "da." The brain "fuses" the two inputs into a third plausible result. While not directly relevant to VR, this effect illustrates how sensory fusion works.
\end{itemize}

\subsubsection{Loudness Perception}

\textit{Loudness perception} varies with frequency. This must be considered in audio interaction experiments, referencing \textbf{equal-loudness contours} (also known as \textit{Fletcher-Munson curves}), which show how perceived loudness changes across different frequencies at the same \textit{sound pressure level}.

\subsubsection{Steven's Power Law}

The response to a stimulus (\textit{perceived magnitude}) is not linear but follows a \textit{power law}:
\begin{equation}
    p = k \cdot m^x
\end{equation}
where $p$ is the perceived magnitude, $m$ is the physical magnitude, $k$ is a constant, and $x$ is an exponent that depends on the stimulus type (e.g., $x = 3.5$ for electric shock, $x = 0.33$ for brightness).

\subsubsection{Just Noticeable Difference (JND)}

The \textbf{Just Noticeable Difference (JND)} is the minimum amount by which a stimulus must be modified for a subject to perceive the change at least 50\% of the time. This threshold varies across sensory modalities and stimulus types.

\subsubsection{Sound Localization}

\textbf{Sound localization} is the estimation of a sound source's position through hearing. This is crucial for VR: auditory and visual stimuli must correspond (e.g., a voice should appear to come from the avatar's mouth).

The JND applied to localization is the \textbf{Minimum Audible Angle (MAA)}, the minimum angular variation that can be detected.

\subsubsection{Monaural cues (one ear)}
\begin{itemize}
    \item \textbf{Pinna (auricle)}: The asymmetric shape of the outer ear distorts incoming sound in a direction-dependent manner, especially for elevation perception.
    
    \item \textbf{Amplitude}: Decreases quadratically with distance.
    
    \item \textbf{Spectral distortion}: At distance, high frequencies attenuate more rapidly (e.g., distant vs. nearby thunder).
    
    \item \textbf{Reverberation}: Can be used for echolocation.
\end{itemize}

\subsubsection{Binaural cues (two ears)}
\begin{itemize}
    \item \textbf{Interaural Level Difference (ILD)}: The difference in sound magnitude (volume) as heard by each ear. The farther ear is in acoustic "shadow" (works better for high frequencies).
    
    \item \textbf{Interaural Time Difference (ITD)}: The difference in arrival time of sound at the two ears (approximately 21.5 cm apart, resulting in a maximum delay of $\sim$0.6 ms).
    
    \item \textbf{Cone of confusion}: A set of positions in space that produce the same ILD and ITD, making localization ambiguous (e.g., a sound in front or behind). Head movement helps resolve this ambiguity.
\end{itemize}

\subsection{Hierarchical Processing}

Stimuli captured by receptors ascend through a \textit{hierarchical network of neurons}. In early stages, signals are combined and propagated upward. In later stages, information flows bidirectionally, allowing \textit{top-down influences} to modulate \textit{lower-level processing}.

\subsection{Information Fusion}

\begin{tcolorbox}[colback=blue!5!white,colframe=blue!75!black,title=\textbf{Multisensory Integration and Sensory Conflict}]
Signals from multiple senses and proprioception are processed and combined by the brain with our experiences. We expect coherent information across modalities. VR often cannot provide all sensations correlated with a simulation, leading to \textbf{sensory conflict}.

This conflict can cause problems such as dizziness or nausea, particularly in cases of \textbf{vection}.
\end{tcolorbox}

\subsubsection{Vection and Optical Flow}

\begin{tcolorbox}[colback=green!5!white,colframe=green!75!black,title=\textbf{Vection: The Illusion of Self-Motion}]
\textbf{Vection} is the illusion of self-motion. A classic example: you are stationary in a train, and the train on the adjacent track moves, giving you the illusion that you are moving backward.

In VR, vision (\textit{optical flow}) tells the brain that you are accelerating, but the \textit{vestibular system} (balance) and \textit{proprioception} indicate that you are stationary. This conflict is a primary cause of \textbf{VR sickness}.

Movement is perceived (also in computer vision) by estimating "optical flow" from image sequences. Flow vectors are used to recover the type of movement (vection). Types of vection/movement include: (a) yaw, (b) pitch, (c) roll, (d) lateral, (e) vertical, (f) forward/backward.
\end{tcolorbox}

\subsubsection{Factors affecting vection sensitivity}
\begin{itemize}
    \item \textbf{Percentage of visual field (FOV)}: The larger the moving area, the greater the vection. A small moving object is perceived as "object in motion," while an entire moving environment is perceived as "I am moving."
    
    \item \textbf{Distance from center of view}: Peripheral motion contributes more to vection perception.
    
    \item \textbf{Exposure time}: Self-motion perception increases with time.
    
    \item \textbf{Spatial frequency}: Greater detail (texture) increases optical flow perception.
    
    \item \textbf{Contrast}: Higher contrast increases the optical flow signal.
    
    \item \textbf{Other sensory cues}: If coherent (e.g., engine noise, vibrations), they can enhance (correct or incorrect) motion perception.
    
    \item \textbf{Prior knowledge}: Knowing what is about to happen can alter perception.
    
    \item \textbf{Attention}: If the user is distracted (e.g., aiming with a weapon, selecting a menu), vection and its side effects can be mitigated.
    
    \item \textbf{Training/prior adaptation}: With sufficient exposure, the body can learn to distinguish vection from real movement, reducing VR sickness.
\end{itemize}

\section{Interaction and Usability}

\subsection{Interaction Design}
%itemize
\begin{itemize}
    \item \textbf{Human Factors}: The study of human capabilities and limitations in the 
    design of human-computer interfaces.
    \item \textbf{Appearance and presentation}: The design of the user interface to make it attractive and easy to use.
    \item \textbf{Interaction}: The design of the interaction between a user and a system.
\end{itemize}    
\subsubsection{Human Factors}
Human factors are the psychological and physiological characteristics of the human body that are relevant to the design of human-computer interfaces.

We need to identify the capabilities and limitations involved in the user activities to solve an interaction task.
In particular we need to focus on the following aspects:
\begin{itemize}
    \item \textbf{Perception}
    \item \textbf{Memory}
    \item \textbf{Attention and Learning}
\end{itemize}
We also need to give control to the user and reduce the amount of \textit{cognitive load}, i.e. the amount of
 information that the user needs to process to solve an \textit{interaction task}.

\subsubsection{User Control}
User control is the ability of the user to control the system.
We need to give the user control and guidance over the system.
The user is the one that decides which functions to use and when to use them.
The interface should differentiate between different users, expertise levels and tasks building a \textbf{flexible interface}.
\subsubsection{Reducing Cognitive Load}
\begin{tcolorbox}[colback=cyan!5!white,colframe=cyan!75!black,title=\textbf{Recognition vs. Recall}]
It is easier for a user to \textit{recognize} rather than \textit{remember}.
It is good practice to show the user all the available commands and options, show proper information to complete a task and organize
the information to help the user focus on the \textit{task at hand}.
\end{tcolorbox}


\subsection{Appearance and Presentation}
We need to distinguish between 3 different principles:
\begin{enumerate}
    \item Create an attractive aspect.
    \item Use meaningful and recognizable object representations.
    \item Keep the interface simple and consistent.
\end{enumerate}

\subsubsection{Create an attractive aspect}
\begin{itemize}
\item All elements on the screen should be visible and recognizable.
\item A proper use of contrast
\item A proper use of colors, shapes and sizes.
\end{itemize}

\subsubsection{Use meaningful and recognizable object representations}
\begin{itemize}
    \item Use clear and distinguishable objects.
    \item Group hierarchically the interface elements.
    \item A visual coherence should be achieved.
\end{itemize}

\begin{tcolorbox}[colback=teal!5!white,colframe=teal!75!black,title=\textbf{Functional Consistency}]
We need to obtain \textit{functional consistency}, i.e. the same function should be represented by the same object in the same way.
A consistent design provides a more \textit{predictable} and \textit{intuitive} interface and thus it is easier to learn for the user.
Icons should suggest how to use the interface. The use of non-consistent icons is a source of \textit{confusion} and \textit{error} for the user.
\end{tcolorbox}

\subsection{Interaction}
It is the way the user interacts with the interface. This can be achieved through the use of mouse, keyboard, touch screen, voice, etc.
\begin{tcolorbox}[colback=orange!5!white,colframe=orange!75!black,title=\textbf{Core Interaction Design Principles}]
We distinguish between 3 main different design principles:
\begin{itemize}
    \item \textbf{Direct handling}: The user can directly handle the interface elements. The aim is to mimic the \textit{natural way} of interacting with the world.
    \item \textbf{Feedback}: The user should receive \textit{feedback} from the interface. This feedback should be \textit{immediate} and \textit{clear}. This also improves 
    the user \textit{awareness} and \textit{learning}. Some examples of feedback are: highlighting a selected object, make a sound when an action is completed, etc.
    \item \textbf{Error tolerance}: The user should be able to recover from \textit{errors}. The system can accept wrong actions without negative consequences.
\end{itemize}
\end{tcolorbox}

There are also other design principles that are important to consider such as:
\begin{itemize}
    \item \textbf{Affordance}: help the user to understand the interface and the \textit{actions} that can be performed.
    \item \textbf{Causality}: help the user to understand the relationship between the actions and the results. If false causality occurs,
    we get an effect called \textit{overshooting}.
    \item \textbf{Constraints}: help the user to understand the \textit{limitations} of the interface.
    \item \textbf{Conceptual model}: 
    \begin{tcolorbox}[colback=violet!5!white,colframe=violet!75!black,title=\textbf{Conceptual Model}]
    Help the user to understand the interface and the actions that can be performed. 
    This is a \textit{mental model} of the interface that the user can use to understand the actions that can be performed.
    A good conceptual model enables the user to predict the results of the actions that can be performed and help him understand
    relations between the device control and outputs.
    A bad conceptual model can lead to confusion and error for the user, it enforces the user to operate blindly and makes
    the identification of action effects more difficult.
    \end{tcolorbox}
    \item \textbf{Visibility}: it is achieved by placing the controls in a \textit{high visibility area}. Just by looking
    at the interface the user can see the controls and the actions. Visibility is often violated in order to make things look more appealing.
    \item \textbf{Mapping}: it is a relation between two objects. It is important that this is easy to 
    learn and remember. Some types of mappings are \textit{spatial} and \textit{perceptual}.
    
\end{itemize}
\subsection{Classic Interaction and pervasive interfaces}
Classic interaction is the traditional way of interacting with the interface. 
Traditional computer interfaces are based on the \textit{Desktop concept}, as an example of 
\textit{remediation}. Together with the Desktop concept we have the icons, mouse, menus, etc.
\subsubsection{Pervasive interfaces}
In pervasive interfaces, the computational infrastructure is distributed in the environment.
The interface identifies and serves the user needs.
Other similar concepts are \textit{ubiquitous computing}, the \textit{invisible computer}, the \textit{disappearing computer}, \textit{ambient intelligence}, and others.
Computation is part of our everyday life, standard objects became \textit{computational objects}.
\begin{itemize}
    \item \textbf{Calm Computing}: the computer does not absorb the user's full attention, but requires an intermittent level of attention.
\end{itemize}
We need to put some limits to this \textbf{Pervasive Computing}, in order to avoid the user to be overwhelmed by the amount of
information and the complexity of the interface. We also need to guarantee the user \textbf{privacy} and \textbf{security}.

\begin{tcolorbox}[colback=purple!5!white,colframe=purple!75!black,title=\textbf{Implicit Interaction in Pervasive Computing}]
In pervasive computing the user interaction with the machine changes:
\begin{itemize}
    \item \textbf{Implicit input}: human \textit{actions} and \textit{behaviors} are captured, recognized and interpreted by the machine as inputs.
    \item \textbf{Implicit output}: the machine generates outputs that are not directly visible but are \textit{integrated in the environment}.
\end{itemize}
\end{tcolorbox}

Some examples of pervasive interfaces are:
\begin{itemize}
    \item \textbf{wearable devices}: smartwatches, smartglasses, smartrings, etc.
    \item \textbf{tangible computing}: computation in everyday objects.
    \item \textbf{Locative media}: media and information that are tied to a \textit{specific location}.
    \item \textbf{Intelligent Environments}: the computer is part of the environment and it is \textit{integrated} in it.
\end{itemize}
\textbf{Natural interfaces} are more effective for the user since the user is focused on the solution of the task at hand rather than understanding the interface.
The idea of pervasive interaction is nowadays feasible with the development of the \textbf{Internet of Things (IoT)} and \textbf{Artificial Intelligence (AI)}.
The HCI solution must exploit \textit{AI} and \textit{machine learning} techniques to understand the user's needs and to provide the user with the best possible experience.


\section{Feature Matching}

Feature matching is the process of finding the best match between two sets of features. In images, this means finding the best matching keypoints between two sets. To correctly understand this chapter, we need to know the concept of \textbf{convolution}.

\subsection{Convolution}

\begin{tcolorbox}[colback=cyan!5!white,colframe=cyan!75!black,title=\textbf{Convolution Operation}]
Convolution is a fundamental mathematical operation used extensively in image processing. It combines two functions to produce a third function that expresses how the shape of one is modified by the other.
\end{tcolorbox}

Convolution is a mathematical operation between two functions defined as follows:

\begin{equation}
    (f * g)(t) = \int_{0}^{t} f(\tau) g(t - \tau) d\tau \quad \text{for } f,g : [0,\infty] \rightarrow \mathbb{R}
\end{equation}

where $f$ and $g$ are the two functions to convolve. The meaning of this operation is that the output is the area under the product of the two functions, where one function is flipped and shifted.

In images, convolution is used to apply a filter to the image. The filter is a small matrix of weights (also called a kernel) that is applied to the image. The filter is applied by sliding it over the image and computing the dot product of the filter and the image patch at each position. The result is a new image with the same size as the original image.

\begin{tcolorbox}[colback=purple!5!white,colframe=purple!75!black,title=\textbf{Common Image Filters}]
Some examples of filters used in image processing:
\begin{itemize}
    \item \textbf{Blurring}: A Gaussian filter that blurs the image by averaging neighboring pixels
    \item \textbf{Sharpening}: A Laplacian filter that enhances edges and fine details
    \item \textbf{Edge detection}: A Sobel filter that detects edges by computing gradient magnitudes
\end{itemize}
\end{tcolorbox}

Convolution in images is crucial to build image filters for the estimation of corresponding points between two images.

\subsection{Corresponding Point Computation}

\begin{tcolorbox}[colback=teal!5!white,colframe=teal!75!black,title=\textbf{Feature Matching Pipeline}]
The idea consists in identifying a set of salient points (feature points) to be used for matching. The process involves three main phases:
\begin{enumerate}
    \item \textbf{Salient point detection}: Identification of distinctive points in the image that can be reliably detected
    \item \textbf{Salient point description}: Creation of a descriptor that uniquely characterizes each detected point
    \item \textbf{Salient point matching}: Comparison and matching of descriptors between two images to find corresponding points
\end{enumerate}
\end{tcolorbox}

\subsubsection{SIFT (Scale-Invariant Feature Transform)}

\begin{tcolorbox}[colback=blue!5!white,colframe=blue!75!black,title=\textbf{SIFT Overview}]
SIFT (Scale-Invariant Feature Transform) is a powerful feature detection and description algorithm used to detect and describe salient points (keypoints) in images. The algorithm is designed to be invariant to:
\begin{itemize}
    \item \textbf{Scale}: Features can be detected regardless of the image scale
    \item \textbf{Rotation}: Features remain stable under image rotation
    \item \textbf{Illumination}: Features are robust to changes in lighting conditions
    \item \textbf{Affine transformation}: Features are stable under affine transformations
\end{itemize}
\end{tcolorbox}

The main objective of SIFT is to detect the \textbf{location} and \textbf{scale} of stable points in the image. These stable points are called \textbf{keypoints} or \textbf{feature points}, and they represent distinctive regions in the image that can be reliably detected and matched across different views.

\subsubsection{Scale Space Representation}

The fundamental concept behind SIFT is the \textbf{scale space}, which is a multi-scale representation of an image. The scale space is created by convolving the input image with Gaussian filters at different scales (standard deviations). 

The scale space function $L(x, y, \sigma)$ is defined as:

\begin{equation}
    L(x, y, \sigma) = G(x, y, \sigma) * I(x, y)
\end{equation}

where:
\begin{itemize}
    \item $G(x, y, \sigma)$ is the Gaussian filter (kernel) with standard deviation $\sigma$:
    \[
        G(x, y, \sigma) = \frac{1}{2\pi\sigma^2} e^{-\frac{x^2 + y^2}{2\sigma^2}}
    \]
    \item $I(x, y)$ is the input image
    \item $*$ denotes the convolution operation
    \item $\sigma$ is the scale parameter (larger $\sigma$ means more blurring)
\end{itemize}

The scale space is essentially a \textbf{stack of images}, where each image in the stack is the original image blurred by a Gaussian filter with a different scale parameter $\sigma$. As $\sigma$ increases, the image becomes more blurred, effectively removing fine details and keeping only the larger structures.

\subsubsection{Image Pyramid and Octaves}

To efficiently compute the scale space, SIFT uses an \textbf{image pyramid} structure. The pyramid is organized into \textbf{octaves}, where each octave represents a doubling of the scale. Within each octave, multiple scale levels are computed by progressively increasing $\sigma$ by a constant factor $k$ (typically $k = \sqrt{2}$).

\begin{tcolorbox}[colback=green!5!white,colframe=green!75!black,title=\textbf{Image Pyramid Structure}]
\begin{itemize}
    \item \textbf{Octave}: A set of images at different scales, where the scale doubles between consecutive octaves
    \item \textbf{Scale levels}: Within each octave, multiple blurred versions are created with increasing $\sigma$ values
    \item \textbf{Downsampling}: After each octave, the image is downsampled by a factor of 2 to create the next octave
    \item This structure allows efficient computation while maintaining scale invariance
\end{itemize}
\end{tcolorbox}

\subsubsection{Difference of Gaussians (DoG)}

Instead of directly computing the Laplacian of Gaussian (LoG), which is computationally expensive, SIFT uses the \textbf{Difference of Gaussians (DoG)} as an efficient approximation. The DoG is computed by subtracting two adjacent scale-space images:

\begin{equation}
    \text{DoG}(x, y, \sigma) = L(x, y, k\sigma) - L(x, y, \sigma)
\end{equation}

This can be rewritten using the convolution property:

\begin{equation}
    \text{DoG}(x, y, \sigma) = (G(x, y, k\sigma) - G(x, y, \sigma)) * I(x, y)
\end{equation}

\subsubsection{Keypoint Detection: Finding Extrema}

The keypoints (salient points) in SIFT are detected as \textbf{local extrema} (maxima or minima) in the DoG scale space. This means we look for points that are either brighter or darker than all their neighbors in both spatial and scale dimensions.

\begin{tcolorbox}[colback=orange!5!white,colframe=orange!75!black,title=\textbf{Extrema Detection Process}]
To detect a keypoint at position $(x, y)$ and scale $\sigma$:
\begin{enumerate}
    \item Compare the pixel value at $(x, y, \sigma)$ in the DoG with its \textbf{26 neighbors}:
    \begin{itemize}
        \item 8 neighbors in the same scale level (spatial neighbors)
        \item 9 neighbors in the scale level above ($\sigma \cdot k$)
        \item 9 neighbors in the scale level below ($\sigma / k$)
    \end{itemize}
    \item A point is a keypoint candidate if it is either:
    \begin{itemize}
        \item A \textbf{local maximum}: greater than all 26 neighbors, or
        \item A \textbf{local minimum}: smaller than all 26 neighbors
    \end{itemize}
    \item This creates a \textbf{3D search space} (2D spatial + 1D scale), ensuring that keypoints are stable across both position and scale
\end{enumerate}
\end{tcolorbox}

The 3D neighborhood check ensures that detected keypoints are:
\begin{itemize}
    \item \textbf{Spatially stable}: The point is distinctive in its local spatial neighborhood
    \item \textbf{Scale stable}: The point persists across different scales, making it robust to scale changes
    \item \textbf{Distinctive}: Only truly salient features are selected, reducing false positives
\end{itemize}

\subsubsection{Stability and Robustness}

The stability of a keypoint is determined by analyzing its behavior in the scale space. Points that remain extrema across multiple scales are considered more stable and reliable. The Hessian matrix of the DoG can be used to further analyze the local structure and filter out unstable keypoints (e.g., those along edges, which are less distinctive than corner-like features).

The key insight is that \textbf{stable keypoints} are those that:
\begin{itemize}
    \item Remain extrema across multiple scale levels
    \item Have sufficient contrast (not too weak)
    \item Are not located on edges (which are less distinctive than corners or blobs)
\end{itemize}

These stable keypoints form the basis for feature matching, as they can be reliably detected and matched across different images, even when the images are taken from different viewpoints, scales, or under different lighting conditions.

\subsubsection{Keypoint Orientation Assignment}

After detecting keypoints, we need to assign an \textbf{orientation} to each keypoint to achieve rotation invariance. This is done by computing the gradient magnitude and direction in a local neighborhood around each keypoint.

The orientation is computed using the Gaussian-smoothed image at the scale where the keypoint was detected (not the DoG image). For each keypoint, we consider a local region (typically 16×16 pixels) around it and compute:

\begin{itemize}
    \item \textbf{Gradient magnitude}: $m(x, y) = \sqrt{(L(x+1, y) - L(x-1, y))^2 + (L(x, y+1) - L(x, y-1))^2}$
    \item \textbf{Gradient direction}: $\theta(x, y) = \arctan((L(x, y+1) - L(x, y-1)) / (L(x+1, y) - L(x-1, y)))$
\end{itemize}

where $L(x, y)$ is the Gaussian-smoothed image at the keypoint's scale.

A \textbf{histogram of gradient orientations} is then created from this local region, with 36 bins covering 360 degrees. The dominant orientation (the peak of the histogram) is assigned to the keypoint. If there are multiple peaks above 80\% of the maximum peak, multiple keypoints are created at the same location with different orientations.

\begin{tcolorbox}[colback=yellow!10!white,colframe=yellow!75!black,title=\textbf{Orientation Assignment}]
\begin{itemize}
    \item Computed from the Gaussian-smoothed image (not DoG) at the keypoint's scale
    \item Uses a 16×16 pixel local region around the keypoint
    \item Creates a 36-bin histogram of gradient orientations
    \item Assigns the dominant orientation (histogram peak) to the keypoint
    \item Multiple orientations can be assigned if there are secondary peaks above 80\% of the maximum
\end{itemize}
\end{tcolorbox}

\subsubsection{SIFT Descriptor}

The \textbf{SIFT descriptor} is a 128-dimensional vector that uniquely describes the local appearance around each keypoint. The descriptor is designed to be robust to illumination changes, small geometric distortions, and viewpoint variations.

The descriptor is computed as follows:

\begin{enumerate}
    \item \textbf{Local region}: A 16×16 pixel region around the keypoint is considered, rotated according to the keypoint's orientation to achieve rotation invariance.
    
    \item \textbf{Subdivision}: This region is divided into 4×4 sub-regions (16 sub-regions total), each containing 4×4 pixels.
    
    \item \textbf{Gradient histogram}: For each 4×4 sub-region, an 8-bin histogram of gradient orientations is computed. The gradient magnitudes are weighted by a Gaussian window centered on the keypoint to give more weight to pixels closer to the keypoint.
    
    \item \textbf{Concatenation}: The 16 histograms (one per sub-region), each with 8 bins, are concatenated to form a 128-dimensional vector (16 × 8 = 128).
    
    \item \textbf{Normalization}: The descriptor vector is normalized to unit length to achieve illumination invariance. To reduce the effect of non-linear illumination changes, values are thresholded (typically at 0.2) and renormalized.
\end{enumerate}

\begin{tcolorbox}[colback=orange!5!white,colframe=orange!75!black,title=\textbf{SIFT Descriptor Properties}]
\begin{itemize}
    \item \textbf{Dimensionality}: 128 elements (16 sub-regions × 8 orientation bins)
    \item \textbf{Rotation invariant}: The region is rotated according to keypoint orientation
    \item \textbf{Illumination robust}: Normalized to unit length and thresholded
    \item \textbf{Geometrically stable}: Uses Gaussian weighting to emphasize central pixels
    \item \textbf{Distinctive}: Captures local gradient patterns that uniquely identify the feature
\end{itemize}
\end{tcolorbox}

\subsubsection{Feature Matching}

Once descriptors are computed for keypoints in both images, matching is performed by computing the distance between descriptor vectors. The most common distance metric is the \textbf{Euclidean distance} (L2 norm) between the 128-dimensional vectors.

For each keypoint in the first image, we find its nearest neighbor (smallest distance) and second nearest neighbor in the second image. A match is accepted if the ratio of 
these distances is below a threshold (typically 0.8). This helps eliminate ambiguous matches and improves matching reliability.

Nowadays, there are many new methods for feature matching obtained from deep learning approaches (such as SuperPoint, LoFtr...) which can provide better performance in certain scenarios.

\section{Feature Matching with Deep Learning}

Traditional feature matching methods like SIFT rely on hand-crafted features and descriptors. Deep learning approaches have revolutionized feature matching by learning optimal feature representations directly from data, often achieving superior performance in challenging scenarios such as low-texture regions, extreme viewpoint changes, and varying illumination conditions.

\begin{tcolorbox}[colback=blue!5!white,colframe=blue!75!black,title=\textbf{Advantages of Deep Learning for Feature Matching}]
\begin{itemize}
    \item \textbf{Learned representations}: Features are learned from data rather than hand-designed
    \item \textbf{Better generalization}: Can handle challenging scenarios better than traditional methods
    \item \textbf{End-to-end training}: Can optimize the entire pipeline jointly
    \item \textbf{Context awareness}: Can leverage semantic and contextual information
    \item \textbf{Robustness}: Often more robust to illumination, viewpoint, and scale changes
\end{itemize}
\end{tcolorbox}

\subsection{SuperPoint}

SuperPoint is a self-supervised framework for training interest point detectors and descriptors. It uses a fully convolutional neural network that operates on full-sized images and simultaneously computes pixel-level interest point locations and associated descriptors in a single forward pass.

\subsubsection{Architecture}

The SuperPoint architecture consists of:
\begin{itemize}
    \item A \textbf{shared encoder} that processes the input image to extract features
    \item A \textbf{detector head} that predicts interest point locations and their scores
    \item A \textbf{descriptor head} that computes dense descriptors at each pixel location
\end{itemize}

The network outputs:
\begin{itemize}
    \item A \textbf{heatmap} indicating the probability of an interest point at each pixel
    \item A \textbf{descriptor map} with 256-dimensional descriptors for each pixel
\end{itemize}

\subsubsection{Self-Supervised Training}

SuperPoint uses a self-supervised training approach:
\begin{enumerate}
    \item \textbf{Synthetic pre-training}: The network is first pre-trained on synthetic geometric shapes (MagicPoint) to learn basic interest point detection
    \item \textbf{Homographic adaptation}: Real images are warped with random homographies, and the network learns to detect points that are consistent across these transformations
    \item \textbf{Descriptor learning}: Descriptors are trained to be similar for corresponding points and dissimilar for non-corresponding points
\end{enumerate}

\begin{tcolorbox}[colback=green!5!white,colframe=green!75!black,title=\textbf{SuperPoint Key Features}]
\begin{itemize}
    \item \textbf{Single forward pass}: Detects keypoints and computes descriptors simultaneously
    \item \textbf{Dense descriptors}: Provides descriptors at every pixel, not just at detected keypoints
    \item \textbf{Real-time capable}: Fast enough for real-time applications
    \item \textbf{Self-supervised}: Does not require manual annotation of correspondences
\end{itemize}
\end{tcolorbox}

\subsection{LoFTR (Loosely-coupled Feature Transform)}

LoFTR (Loosely-coupled Feature Transform) is a detector-free method for local feature matching that revolutionized dense correspondence estimation. Unlike traditional methods that follow a detect-then-describe pipeline (first detect keypoints, then compute descriptors), LoFTR directly establishes correspondences between image pairs without explicit keypoint detection.

\begin{tcolorbox}[colback=blue!5!white,colframe=blue!75!black,title=\textbf{LoFTR Core Concept}]
LoFTR treats feature matching as a \textbf{set-to-set matching problem} rather than a point-to-point matching problem. Instead of detecting sparse keypoints, it operates on dense feature maps and uses transformer attention to establish correspondences, making it particularly effective in low-texture regions where traditional methods struggle.
\end{tcolorbox}

\subsubsection{Key Innovation}

LoFTR's main innovation lies in the combination of:
\begin{itemize}
    \item \textbf{Transformer architecture}: Uses self-attention and cross-attention mechanisms to model long-range dependencies and establish correspondences
    \item \textbf{Coarse-to-fine strategy}: First matches at low resolution (computationally efficient), then refines at higher resolution (high precision)
    \item \textbf{Detector-free approach}: Eliminates the need for explicit keypoint detection, enabling dense matching even in textureless regions
    \item \textbf{Positional encoding}: Incorporates geometric information directly into the feature representation
\end{itemize}

\subsubsection{Architecture Overview}

The LoFTR pipeline consists of four main stages:

\begin{enumerate}
    \item \textbf{Local feature extraction}
    \item \textbf{Positional encoding}
    \item \textbf{LoFTR module (transformer)}
    \item \textbf{Coarse-to-fine matching}
\end{enumerate}

\subsubsection{Stage 1: Local Feature Extraction}

A CNN backbone (typically ResNet or FPN - Feature Pyramid Network) extracts dense feature maps from both input images. The backbone produces feature maps at multiple scales:
\begin{itemize}
    \item \textbf{Coarse level}: Low-resolution feature maps (typically 1/8 of the original image size) for initial matching
    \item \textbf{Fine level}: High-resolution feature maps (typically 1/2 of the original image size) for refinement
\end{itemize}

The feature maps are flattened into sequences of feature vectors, where each vector corresponds to a spatial location in the image.

\subsubsection{Stage 2: Positional Encoding}

Since transformers are permutation-invariant, positional information must be explicitly added. 
LoFTR uses a \textbf{2D sinusoidal positional encoding} that encodes both the row and column positions of each feature.
This encoding allows the transformer to understand the spatial relationships between features.

\subsubsection{Stage 3: LoFTR Module (Transformer)}

The core of LoFTR is a transformer-based module that processes the feature sequences from both images. The module consists of multiple layers, each containing:

\begin{itemize}
    \item \textbf{Self-attention}: Each image attends to its own features, allowing the model to understand local context and relationships within each image
    \item \textbf{Cross-attention}: Features from one image attend to features from the other image, establishing correspondences between the two images
    \item \textbf{Feed-forward network}: A standard MLP that processes the attended features
\end{itemize}

The attention mechanism computes:
\[
\text{Attention}(Q, K, V) = \text{softmax}\left(\frac{QK^T}{\sqrt{d_k}}\right)V
\]

where:
\begin{itemize}
    \item $Q$ (Query), $K$ (Key), $V$ (Value) are learned linear projections of the input features
    \item $d_k$ is the dimension of the key vectors
    \item The softmax normalizes attention weights
\end{itemize}

The transformer iteratively refines the features through multiple layers, progressively improving the correspondence quality.

\subsubsection{Stage 4: Coarse-to-Fine Matching}

LoFTR uses a two-stage matching strategy:

\textbf{Coarse Matching:}
\begin{itemize}
    \item Operates on low-resolution feature maps (1/8 scale)
    \item Computes a \textbf{similarity matrix} between all feature pairs from the two images
    \item Uses \textbf{dual-softmax} operation to compute matching probabilities:
    \[
    P(i \leftrightarrow j) = \text{softmax}(S_{i,\cdot})_j \cdot \text{softmax}(S_{\cdot,j})_i
    \]
    where $S$ is the similarity matrix and $P(i \leftrightarrow j)$ is the probability that location $i$ in image 1 matches location $j$ in image 2
    \item Selects matches with high confidence (mutual nearest neighbors)
\end{itemize}

\textbf{Fine Matching:}
\begin{itemize}
    \item For each coarse match, extracts a local patch from the high-resolution feature maps
    \item Applies the LoFTR module again on these local patches to refine the correspondence
    \item Achieves sub-pixel accuracy by interpolating the matching location
\end{itemize}

\begin{tcolorbox}[colback=green!5!white,colframe=green!75!black,title=\textbf{Coarse-to-Fine Strategy Benefits}]
\begin{itemize}
    \item \textbf{Efficiency}: Coarse matching reduces the search space from $H \times W \times H \times W$ to a much smaller set of candidate matches
    \item \textbf{Accuracy}: Fine matching refines correspondences to sub-pixel precision
    \item \textbf{Robustness}: Coarse matching provides good initialization even in challenging scenarios
    \item \textbf{Scalability}: Can handle high-resolution images efficiently
\end{itemize}
\end{tcolorbox}

\subsubsection{Training}

LoFTR is trained in a supervised manner using ground-truth correspondences from datasets like MegaDepth or ScanNet. The training loss combines:
\begin{itemize}
    \item \textbf{Matching loss}: Encourages correct correspondences
    \item \textbf{Confidence loss}: Ensures high confidence for correct matches and low confidence for incorrect ones
\end{itemize}

The model learns to:
\begin{itemize}
    \item Extract discriminative features that are robust to viewpoint and illumination changes
    \item Use attention to focus on relevant regions for matching
    \item Establish correspondences even in challenging scenarios (low texture, occlusions, etc.)
\end{itemize}

\subsubsection{Advantages}

\begin{tcolorbox}[colback=orange!5!white,colframe=orange!75!black,title=\textbf{LoFTR Advantages}]
\begin{itemize}
    \item \textbf{Detector-free}: No need for explicit keypoint detection, enabling matching in textureless regions
    \item \textbf{Dense matching}: Can establish correspondences at every pixel location, not just at detected keypoints
    \item \textbf{Global context}: Transformer attention mechanism captures long-range dependencies and geometric relationships
    \item \textbf{Robust to viewpoint changes}: Handles large viewpoint differences (up to 60° rotation) better than local methods
    \item \textbf{Sub-pixel accuracy}: Achieves sub-pixel matching precision through fine-level refinement
    \item \textbf{Handles low-texture regions}: Particularly effective in areas with repetitive patterns or uniform textures
    \item \textbf{End-to-end trainable}: The entire pipeline can be optimized jointly
\end{itemize}
\end{tcolorbox}


Despite its strengths, LoFTR has some limitations:
\begin{itemize}
    \item \textbf{Computational cost}: Transformer attention is computationally expensive, especially for high-resolution images
    \item \textbf{Memory requirements}: Dense feature maps require significant memory
    \item \textbf{Training data}: Requires large amounts of training data with ground-truth correspondences
    \item \textbf{Real-time performance}: May be slower than traditional methods for real-time applications
\end{itemize}

\subsubsection{Applications}

LoFTR has been successfully applied to:
\begin{itemize}
    \item \textbf{Visual localization}: Camera pose estimation from images
    \item \textbf{Structure from Motion}: 3D reconstruction from image sequences
    \item \textbf{Image matching}: Finding correspondences for image stitching and panorama creation
    \item \textbf{SLAM}: Visual SLAM systems that require robust feature matching
    \item \textbf{Augmented Reality}: Camera tracking and scene understanding
\end{itemize}


\subsection{Comparison with Traditional Methods}

\begin{tcolorbox}[colback=purple!5!white,colframe=purple!75!black,title=\textbf{Deep Learning vs. Traditional Methods}]
\begin{center}
\begin{tabular}{|l|l|l|}
\hline
\textbf{Aspect} & \textbf{Traditional (SIFT)} & \textbf{Deep Learning} \\
\hline
Feature design & Hand-crafted & Learned from data \\
\hline
Keypoint detection & Explicit (DoG extrema) & Implicit or learned \\
\hline
Descriptor & Fixed 128-dim vector & Learned, variable size \\
\hline
Training & No training needed & Requires training data \\
\hline
Generalization & Limited & Better on diverse data \\
\hline
Low-texture regions & Struggles & Often better \\
\hline
Computational cost & Moderate & Can be higher \\
\hline
Real-time capability & Yes & Depends on architecture \\
\hline
\end{tabular}
\end{center}
\end{tcolorbox}

\subsection{Applications}

Deep learning-based feature matching has found applications in:
\begin{itemize}
    \item \textbf{Visual SLAM}: Simultaneous Localization and Mapping using visual features
    \item \textbf{Structure from Motion}: 3D reconstruction from image sequences
    \item \textbf{Image stitching}: Panorama creation
    \item \textbf{Augmented Reality}: Camera pose estimation and tracking
    \item \textbf{Visual odometry}: Motion estimation from camera images
    \item \textbf{Object tracking}: Tracking objects across frames
\end{itemize}

\section{Fundamentals of Computer Vision}

Computer vision is the field of study that focuses on enabling computers to interpret and understand visual information from the world. It involves extracting meaningful information from images and videos, which is fundamental for many HCI applications, including augmented reality, gesture recognition, and visual tracking.

\subsection{Image Acquisition and Representation}

This section introduces how the physical world is converted into digital data. Understanding the image acquisition process is fundamental to computer vision.

\subsubsection{Digital Images}

Digital images are defined as \textbf{2D arrays} (matrices) of numbers that represent the light intensity recorded by a photosensitive device. Each element in the array is called a \textbf{pixel} (picture element).

\begin{tcolorbox}[colback=blue!5!white,colframe=blue!75!black,title=\textbf{Digital Image Representation}]
\begin{itemize}
    \item Images are stored as 2D matrices $I(x, y)$ where $(x, y)$ are pixel coordinates
    \item Each pixel value represents light intensity (brightness) at that location
    \item For grayscale images: single value per pixel (typically 0-255 for 8-bit images)
    \item For color images: multiple values per pixel (e.g., RGB with red, green, blue channels)
\end{itemize}
\end{tcolorbox}

\subsubsection{Optical Process}

The optical process describes the path of light rays through the camera system:

\begin{itemize}
    \item \textbf{Lenses}: Light rays from the 3D world pass through camera lenses, which focus the light
    \item \textbf{Aperture (pupilla)}: The opening that controls how much light enters the camera
    \item \textbf{Image plane}: The surface where the image is formed
    \item \textbf{Color filters}: Modern cameras use \textbf{Bayer pattern} filters on sensors to capture color information. The Bayer pattern arranges red, green, and blue filters in a specific mosaic pattern
    \item \textbf{Infrared filters}: Many cameras include IR filters to block infrared light and improve color accuracy
\end{itemize}

\subsubsection{Focus and Aperture}

The concept of \textbf{focus} is crucial for image formation:

\begin{tcolorbox}[colback=green!5!white,colframe=green!75!black,title=\textbf{Point in Focus}]
For a point to be in focus, all light rays from a 3D point $\mathbf{P}$ must converge to a single image point $\mathbf{p}$ on the image plane.
\end{tcolorbox}

\textbf{Aperture and Sharpness:}
\begin{itemize}
    \item By reducing the aperture to a point (pinhole), we obtain a one-to-one correspondence between visible points and image points
    \item This increases sharpness (all points are in focus) but reduces the amount of light entering the camera
    \item Larger apertures allow more light but create depth-of-field effects (only points at a specific distance are in focus)
\end{itemize}

\subsubsection{Sensors}

The image is discretized using photosensitive sensors:

\begin{itemize}
    \item \textbf{CCD (Charge-Coupled Device)}: Converts light energy into electrical charge, which is then read out and converted to digital values
    \item \textbf{CMOS (Complementary Metal-Oxide-Semiconductor)}: Each pixel has its own amplifier, allowing faster readout and lower power consumption
\end{itemize}

The conversion process:
\begin{enumerate}
    \item Light photons hit the sensor, generating electrical charge proportional to light intensity
    \item The charge is converted to voltage
    \item The voltage is quantized into discrete integer values (e.g., 0-255 for 8-bit images)
    \item These values form the digital image matrix
\end{enumerate}

\subsection{Image Geometry (Pinhole Model)}

This section defines the mathematical model that relates 3D world coordinates to 2D image coordinates.

\subsubsection{Camera Obscura}

The simplest geometric model is the \textbf{pinhole camera}, derived from the Renaissance principle of the \textit{camera obscura}. In this model:

\begin{itemize}
    \item The camera has an infinitesimally small aperture (pinhole)
    \item All light rays pass through a single point (the optical center)
    \item The image is formed on a plane behind the pinhole
    \item There is a one-to-one correspondence between points in the 3D world and points on the image plane
\end{itemize}

\subsubsection{Perspective Projection Equations}

Given a 3D point $\mathbf{M} = [X, Y, Z]^T$ in the world coordinate system and its projection $\mathbf{m} = [u, v]^T$ on the image plane, the perspective projection equations are:

\[
u = -f \frac{X}{Z}, \quad v = -f \frac{Y}{Z}
\]

where:
\begin{itemize}
    \item $f$ represents the \textbf{focal length} (distance between the center of projection and the image plane)
    \item The negative sign indicates that the image is inverted (or we can place the image plane in front of the pinhole to avoid inversion)
    \item These relations are \textbf{non-linear} due to the division by the depth $Z$
\end{itemize}

\begin{tcolorbox}[colback=orange!5!white,colframe=orange!75!black,title=\textbf{Perspective Projection Properties}]
\begin{itemize}
    \item Points farther from the camera appear smaller (perspective foreshortening)
    \item Parallel lines in the world converge to a vanishing point in the image
    \item The projection depends on the depth $Z$, making it a non-linear transformation
\end{itemize}
\end{tcolorbox}

\subsubsection{Homogeneous Coordinates}

By introducing \textbf{homogeneous coordinates}, the projection becomes a linear operation. A 3D point $\mathbf{M} = [X, Y, Z]^T$ is represented in homogeneous coordinates as $\mathbf{M} = [X, Y, Z, 1]^T$.

The projection can then be written as:

\[
\mathbf{m} \simeq \mathbf{P} \mathbf{M}
\]

where:
\begin{itemize}
    \item $\mathbf{m} = [u, v, w]^T$ is the projection in homogeneous coordinates
    \item $\mathbf{P}$ is the 3×4 projection matrix
    \item $\simeq$ means equality up to a scale factor
    \item The actual image coordinates are: $u_{img} = u/w$, $v_{img} = v/w$
\end{itemize}

In matrix form:

\[
\begin{pmatrix} u \\ v \\ w \end{pmatrix} = \begin{pmatrix}
p_{11} & p_{12} & p_{13} & p_{14} \\
p_{21} & p_{22} & p_{23} & p_{24} \\
p_{31} & p_{32} & p_{33} & p_{34}
\end{pmatrix} \begin{pmatrix} X \\ Y \\ Z \\ 1 \end{pmatrix}
\]

\subsubsection{Perspective vs Orthographic}

When the depth variation of objects is negligible compared to their distance from the camera ($\Delta Z / Z_0 \ll 1$), we can approximate the perspective projection with a \textbf{scaled orthographic projection} (also called \textbf{weak perspective}).

In weak perspective:
\[
u = -s \frac{X}{Z_0}, \quad v = -s \frac{Y}{Z_0}
\]

where $s = f/Z_0$ is a constant scale factor, and $Z_0$ is the average depth of the object.

\begin{tcolorbox}[colback=purple!5!white,colframe=purple!75!black,title=\textbf{Weak Perspective Approximation}]
Weak perspective is valid when:
\begin{itemize}
    \item Objects are far from the camera relative to their size
    \item The depth variation within the object is small
    \item This approximation simplifies many computer vision algorithms
\end{itemize}
\end{tcolorbox}

\subsection{Camera Parameters}

The complete mathematical model is decomposed into intrinsic and extrinsic parameters.

\subsubsection{Intrinsic Parameters (Matrix K)}

Intrinsic parameters (encoded in the matrix $\mathbf{K}$) describe the internal properties of the camera and are independent of its position and orientation.

\textbf{Transformation from sensor coordinates (meters) to pixel coordinates:}

The intrinsic parameters include:

\begin{itemize}
    \item \textbf{Focal length in pixels} $(\alpha_u, \alpha_v)$: The focal length expressed in pixel units. If the physical focal length is $f$ (in meters) and the pixel size is $(s_x, s_y)$ (in meters per pixel), then:
    \[
    \alpha_u = \frac{f}{s_x}, \quad \alpha_v = \frac{f}{s_y}
    \]
    \item \textbf{Principal point} $(u_0, v_0)$: The optical center expressed in pixel coordinates (the point where the optical axis intersects the image plane)
    \item \textbf{Skew factor} $s$: Accounts for non-rectangular pixels (usually 0 for modern cameras)
\end{itemize}

The intrinsic matrix $\mathbf{K}$ is:

\[
\mathbf{K} = \begin{pmatrix}
\alpha_u & s & u_0 \\
0 & \alpha_v & v_0 \\
0 & 0 & 1
\end{pmatrix}
\]

\begin{tcolorbox}[colback=cyan!5!white,colframe=cyan!75!black,title=\textbf{Intrinsic Parameters}]
\begin{itemize}
    \item Intrinsic parameters are fixed for a given camera (unless the zoom changes)
    \item They transform points from the camera coordinate system to pixel coordinates
    \item They account for the physical properties of the sensor and lens
\end{itemize}
\end{tcolorbox}

\subsubsection{Extrinsic Parameters (R, t)}

Extrinsic parameters describe the camera's pose (position and orientation) in the world coordinate system.

\begin{itemize}
    \item \textbf{Rotation matrix} $\mathbf{R}$: A 3×3 orthogonal matrix representing the camera's orientation. It transforms coordinates from the world frame to the camera frame
    \item \textbf{Translation vector} $\mathbf{t}$: A 3×1 vector representing the camera's position (the world origin expressed in camera coordinates)
\end{itemize}

The transformation from world coordinates to camera coordinates is:

\[
\mathbf{M}_{cam} = \mathbf{R} \mathbf{M}_{world} + \mathbf{t}
\]

In homogeneous coordinates:

\[
\begin{pmatrix} X_{cam} \\ Y_{cam} \\ Z_{cam} \\ 1 \end{pmatrix} = \begin{pmatrix}
\mathbf{R} & \mathbf{t} \\
\mathbf{0}^T & 1
\end{pmatrix} \begin{pmatrix} X_{world} \\ Y_{world} \\ Z_{world} \\ 1 \end{pmatrix}
\]

\subsubsection{Full Projection Matrix (P)}

The complete projection matrix $\mathbf{P}$ combines both intrinsic and extrinsic parameters:

\[
\mathbf{P} = \mathbf{K} [\mathbf{R} | \mathbf{t}]
\]

where $[\mathbf{R} | \mathbf{t}]$ is a 3×4 matrix formed by concatenating the rotation matrix and translation vector.

The full projection equation is:

\[
\mathbf{m} \simeq \mathbf{P} \mathbf{M}_{world} = \mathbf{K} [\mathbf{R} | \mathbf{t}] \mathbf{M}_{world}
\]

This matrix has 11 degrees of freedom (12 elements minus 1 scale factor) and encodes all information needed to project 3D world points to 2D image coordinates.

\subsection{Camera Calibration (Direct Method)}

This section explains how to estimate the projection matrix $\mathbf{P}$ from known correspondences between 3D points and 2D image points.

\subsubsection{Problem Formulation}

Given $n$ control points $\mathbf{M}_i$ in 3D space and their corresponding projections $\mathbf{m}_i$ in the image, we seek to find the projection matrix $\mathbf{P}$ such that:

\[
\mathbf{m}_i \simeq \mathbf{P} \mathbf{M}_i \quad \text{for } i = 1, \ldots, n
\]

Since this is an equality up to scale, we can write:

\[
\mathbf{m}_i \times (\mathbf{P} \mathbf{M}_i) = \mathbf{0}
\]

where $\times$ denotes the cross product. This eliminates the scale factor and gives us linear constraints.

\subsubsection{Algebraic Tools}

\textbf{Cross Product Formulation:}

The cross product $\mathbf{m}_i \times (\mathbf{P} \mathbf{M}_i) = \mathbf{0}$ gives us two independent linear equations (the third is a linear combination of the first two).

For each correspondence, we can write:

\[
\begin{pmatrix}
0 & -w_i & v_i \\
w_i & 0 & -u_i \\
-v_i & u_i & 0
\end{pmatrix} \mathbf{P} \mathbf{M}_i = \mathbf{0}
\]

where $\mathbf{m}_i = [u_i, v_i, w_i]^T$.

\textbf{Kronecker Product and Vec-Trick:}

To transform the matrix equation into a linear system, we use the \textbf{vec-trick} (vectorization) and \textbf{Kronecker product}:

\[
\text{vec}(\mathbf{A} \mathbf{B} \mathbf{C}) = (\mathbf{C}^T \otimes \mathbf{A}) \text{vec}(\mathbf{B})
\]

where $\otimes$ denotes the Kronecker product and $\text{vec}(\cdot)$ vectorizes a matrix by stacking its columns.

Applying this to our problem, we obtain a linear homogeneous system:

\[
\mathbf{A} \cdot \text{vec}(\mathbf{P}) = \mathbf{0}
\]

where $\mathbf{A}$ is a $2n \times 12$ matrix built from the correspondences, and $\text{vec}(\mathbf{P})$ is the 12-element vector formed by stacking the columns of $\mathbf{P}$.

\subsubsection{Solution}

The system $\mathbf{A} \cdot \text{vec}(\mathbf{P}) = \mathbf{0}$ is solved using \textbf{Singular Value Decomposition (SVD)}:

\begin{enumerate}
    \item Compute the SVD of $\mathbf{A}$: $\mathbf{A} = \mathbf{U} \mathbf{\Sigma} \mathbf{V}^T$
    \item The solution is the right singular vector corresponding to the smallest singular value
    \item This minimizes the least squares error: $\min \|\mathbf{A} \cdot \text{vec}(\mathbf{P})\|^2$
    \item Reshape the solution vector back into a 3×4 matrix $\mathbf{P}$
\end{enumerate}

\begin{tcolorbox}[colback=yellow!10!white,colframe=yellow!75!black,title=\textbf{Calibration Requirements}]
\begin{itemize}
    \item Need at least 6 correspondences (each gives 2 equations, need 11 parameters)
    \item More correspondences improve accuracy and robustness to noise
    \item Points should be well-distributed in 3D space and image
    \item The 3D points should not all lie on a plane (non-coplanar configuration)
\end{itemize}
\end{tcolorbox}

\subsubsection{Calibration Objects}

Practical calibration uses objects with known 3D geometry:

\begin{itemize}
    \item \textbf{Checkerboards}: Planar patterns with known square sizes. Easy to detect corners automatically
    \item \textbf{3D calibration objects}: Objects with known 3D coordinates of markers or corners
    \item \textbf{Multiple views}: Capturing the same calibration pattern from different viewpoints increases the number of constraints
\end{itemize}

\subsection{Stereopsis and Triangulation}

This section analyzes how to reconstruct 3D structure by observing the scene from two different viewpoints.

\subsubsection{Stereopsi}

\textbf{Stereopsis} is the process of estimating 3D structure based on \textbf{binocular disparity} (the difference in position of the same point in the two images).

\begin{tcolorbox}[colback=green!5!white,colframe=green!75!black,title=\textbf{Stereopsis Principle}]
\begin{itemize}
    \item The same 3D point appears at different positions in the left and right images
    \item This horizontal difference is called \textbf{disparity}
    \item Disparity is inversely proportional to depth
    \item By measuring disparity, we can compute the 3D position of points
\end{itemize}
\end{tcolorbox}

\subsubsection{Main Problems}

Stereopsis involves two main problems:

\begin{enumerate}
    \item \textbf{Correspondence (Matching)}: Finding corresponding points between the two images. This is challenging because:
    \begin{itemize}
        \item The same point may look different due to viewpoint changes
        \item Occlusions: some points visible in one image may be hidden in the other
        \item Ambiguity: similar-looking points can be confused
    \end{itemize}
    \item \textbf{Triangulation}: Once correspondences are found, reconstructing the 3D position geometrically
\end{enumerate}

\subsubsection{Triangulation (Parallel Cameras)}

For the simplified case with \textbf{parallel cameras} (aligned cameras with parallel optical axes):

Given a 3D point $\mathbf{M} = [X, Y, Z]^T$:
\begin{itemize}
    \item Left camera sees it at pixel $u$
    \item Right camera sees it at pixel $u'$
    \item The disparity is: $d = u' - u$
\end{itemize}

The depth $Z$ is inversely proportional to the disparity:

\[
Z = \frac{b \cdot f}{u' - u} = \frac{b \cdot f}{d}
\]

where:
\begin{itemize}
    \item $b$ is the \textbf{baseline} (distance between the two camera centers)
    \item $f$ is the focal length (assumed equal for both cameras)
    \item $d = u' - u$ is the disparity
\end{itemize}

Once $Z$ is known, the $X$ and $Y$ coordinates can be computed from either image:

\[
X = \frac{u \cdot Z}{f}, \quad Y = \frac{v \cdot Z}{f}
\]

\subsubsection{Triangulation (General Case - Linear Eigen)}

For the general case with two arbitrary cameras (with known projection matrices $\mathbf{P}$ and $\mathbf{P}'$), a 3D point $\mathbf{M}$ is estimated by solving an overdetermined linear system derived from the projection equations of the two views.

Given:
\begin{itemize}
    \item Point $\mathbf{m} = [u, v]^T$ in the first image (from camera with matrix $\mathbf{P}$)
    \item Point $\mathbf{m}' = [u', v']^T$ in the second image (from camera with matrix $\mathbf{P}'$)
\end{itemize}

The projection equations are:
\[
\mathbf{m} \simeq \mathbf{P} \mathbf{M}, \quad \mathbf{m}' \simeq \mathbf{P}' \mathbf{M}
\]

Using the cross product to eliminate the scale factor:

\[
\mathbf{m} \times (\mathbf{P} \mathbf{M}) = \mathbf{0}, \quad \mathbf{m}' \times (\mathbf{P}' \mathbf{M}) = \mathbf{0}
\]

This gives us 4 linear equations (2 from each view) in 4 unknowns (X, Y, Z, 1). We can write this as:

\[
\mathbf{A} \mathbf{M} = \mathbf{0}
\]

where $\mathbf{A}$ is a 4×4 matrix built from the projection equations.

\textbf{Least Squares Strategy:}

To handle noise in measurements, we use a least squares approach. The solution is found by:

\begin{enumerate}
    \item Building the system $\mathbf{A} \mathbf{M} = \mathbf{0}$ from both views
    \item Computing the SVD: $\mathbf{A} = \mathbf{U} \mathbf{\Sigma} \mathbf{V}^T$
    \item The 3D point $\mathbf{M}$ is the right singular vector corresponding to the smallest singular value
    \item This minimizes the reprojection error in a least squares sense
\end{enumerate}

\begin{tcolorbox}[colback=red!5!white,colframe=red!75!black,title=\textbf{Triangulation Considerations}]
\begin{itemize}
    \item \textbf{Baseline}: Larger baseline improves depth accuracy but makes correspondence harder
    \item \textbf{Viewing angle}: Points should be visible in both views with sufficient angular separation
    \item \textbf{Noise}: Measurement noise affects the accuracy of the reconstructed 3D point
    \item \textbf{Geometric constraints}: Epipolar geometry can help validate correspondences before triangulation
\end{itemize}
\end{tcolorbox}

\section{Augmented Reality}
\textbf{\textit{Mixed Reality}} describes the merging of a real-world environment with a computer-generated one.
Physical and virtual objects coexist and interact in real time, enabling scenarios such as previewing
how a new piece of furniture would look in a room before purchasing it.

\begin{tcolorbox}[colback=blue!4!white,colframe=blue!60!black,title=\textbf{Key Concept: Mixed Reality}]
Successful AR experiences align digital content with the physical world so that
virtual objects respond consistently to real-world geometry, lighting, and user motion.
\end{tcolorbox}

\begin{itemize}
    \item Synthetic and real objects are visualized within the same scenario.
    \item Synthetic objects are rendered through a virtual camera controlled by the graphics pipeline.
    \item Real objects are captured through a physical camera, then processed so they can be displayed within the shared scenario.
\end{itemize}

Integration between the virtual and real scenes is obtained by calibrating the synthetic camera with the same parameters as the physical one.
In practice, both cameras must share the same \textbf{\textit{intrinsic}} and \textbf{\textit{extrinsic}} characteristics.

\subsection{Custom Camera in Unity}
Unity is a game engine that allows the creation of 3D games and simulations (render engine).
We can create a custom camera in Unity to be used in an AR application, modifying its parameters so they match those of the real camera.
Unity exposes both \textbf{\textit{intrinsic}} and \textbf{\textit{extrinsic}} parameters of the camera, which we derive using calibration techniques.

Unity is specialized in video games and simulations, it is not specialized in AR/VR.
Default game objects are:
\begin{itemize}
    \item Object(s)
    \item Camera
    \item Light
\end{itemize}

We particularly focus on the \textbf{\textit{Transform}} component of game objects.
The Transform component positions, rotates, and scales each object in the scene.
The Camera component exposes many tunable parameters that we can modify to match the real camera.

\subsection{Calibration Inputs}
\begin{itemize}
    \item A picture of the real scene.
    \item The intrinsic and extrinsic parameters of the real camera.
    \item A synthetic model represented in the same reference system as the real scene (e.g., a 3D model aligned to the captured environment).
\end{itemize}

We use software such as Zephyr to estimate the camera parameters from the real scene and the synthetic model.

\subsubsection{Calibration Objective}
The virtual camera in the synthetic scene must be configured with the parameters of the real camera.
We therefore control the \textbf{\textit{intrinsic}} and \textbf{\textit{extrinsic}} parameters of the virtual camera to match the real one.
In Unity we mainly need the following \textbf{\textit{intrinsic}} values to control the virtual camera:
\begin{itemize}
    \item Focal length in mm (not in pixels)
    \item Sensor size (pixel size)
    \item Lens shift (principal point) 
\end{itemize}

We apply the correct formulas to convert focal length values from pixels to millimeters.
In Unity we also set the \textbf{\textit{extrinsic}} parameters through the Transform component (position and rotation of the virtual camera).
This step is non-trivial because Zephyr and Unity use different reference systems for the extrinsic parameters.
\begin{itemize}
    \item The x axis is the same in Zephyr and Unity.
    \item The y axis has opposite direction in Zephyr and Unity.
\end{itemize}

We reverse the Y axis using the following matrix:

\begin{equation}
    \begin{pmatrix}
        1 & 0 & 0 \\
        0 & -1 & 0 \\
        0 & 0 & 1
    \end{pmatrix}
\end{equation}

We switch the Z and Y axes using the following matrix:

\begin{equation}
    \begin{pmatrix}
        1 & 0 & 0 \\
        0 & 0 & 1 \\
        0 & 1 & 0
    \end{pmatrix}
\end{equation}

Given the rotation matrix and the translation vector we can compute the extrinsic parameters of the virtual camera.
Finally, in Unity the rotation matrix is converted into Euler angles using a dedicated procedure.

After completing these operations we can convert the Zephyr parameters to Unity parameters (Transform component).

Now we overlap the real scene and the synthetic object by creating a canvas as a child of the camera object.

\begin{tcolorbox}[colback=yellow!8!white,colframe=orange!60!black,title=\textbf{Camera Alignment Checklist}]
\begin{itemize}
    \item Match intrinsic parameters (focal length, sensor size, principal point).
    \item Align extrinsic parameters by mirroring or swapping axes when required.
    \item Validate the pose by rendering a reference object over the captured frame.
\end{itemize}
\end{tcolorbox}

\subsection{Our Pipeline}

\begin{tcolorbox}[colback=green!6!white,colframe=green!50!black,title=\textbf{AR Integration Pipeline}]
\begin{enumerate}
    \item Create or import the 3D representation of the synthetic object.
    \item Calibrate the real camera (e.g., with Zephyr) to obtain intrinsic and extrinsic parameters.
    \item Export the calibration data (e.g., XMP file) and convert focal length values from pixels to millimetres.
    \item Translate the extrinsic parameters into Unity's Transform conventions (Euler angles and axis adjustments).
    \item Render the 3D model with the calibrated virtual camera and composite it with the captured scene.
\end{enumerate}
\end{tcolorbox}
\end{document}  